\documentclass[12pt,a4paper]{report}

\usepackage[utf8]{inputenc}
\usepackage[T1]{fontenc}
\usepackage[italian]{babel}
\usepackage{amsmath}
\usepackage{amssymb}
\usepackage{amsthm}
\usepackage{mathtools}
\usepackage{booktabs}
\usepackage{longtable}
\usepackage{csquotes}
\usepackage{algorithm}
\usepackage{algpseudocode}
\usepackage[hidelinks]{hyperref}

\DeclareMathOperator{\Bernoulli}{Bernoulli}
\DeclareMathOperator{\Uniform}{Uniform}
\DeclareMathOperator{\Cat}{Cat}
\DeclareMathOperator{\MMD}{MMD}
\DeclareMathOperator*{\argmax}{arg\,max}
\DeclareMathOperator*{\argmin}{arg\,min}
\newcommand{\Graphs}{\mathcal{G}}
\newcommand{\GraphDist}{\mathbb{P}}
\newcommand{\Real}{\mathbb{R}}
\newcommand{\KL}{D_{\mathrm{KL}}}
\newcommand{\one}{\mathbf{1}}

\theoremstyle{definition}
\newtheorem{definition}{Definizione}[chapter]
\newtheorem{proposition}[definition]{Proposizione}
\newtheorem{remark}[definition]{Osservazione}

\title{Generazione e valutazione di grafi mediante stima e diffusione di graphon}
\author{TODO: Nome Cognome}
\date{\today}

\begin{document}

\maketitle
\tableofcontents

\chapter{Introduzione}

La generazione di grafi e' un problema centrale nello studio di dati relazionali:
molte applicazioni, dalle reti sociali ai sistemi biologici, richiedono modelli in
grado di produrre nuove istanze che preservino proprieta' strutturali osservate in
un insieme di esempi. A differenza dei dati vettoriali o delle immagini, i grafi
combinano una struttura combinatoria discreta con vincoli di invarianza rispetto
alla permutazione dei nodi. Questa natura rende non banale sia la definizione di
modelli generativi efficaci, sia la loro valutazione quantitativa.

Il presente lavoro considera una pipeline sperimentale per la generazione di grafi
che introduce una \emph{rappresentazione intermedia} basata su graphon. Un graphon
puo' essere interpretato come una funzione simmetrica \(W : [0,1]^2 \to [0,1]\),
dalla quale e' possibile campionare grafi scegliendo variabili latenti associate ai
nodi e generando archi indipendenti con probabilita' data dal valore di \(W\). In
questa prospettiva, il graphon fornisce una rappresentazione continua e non
parametrica della struttura probabilistica di una famiglia di grafi.

L'idea alla base della pipeline studiata in questa tesi, indicata nel repository
come \textbf{OGAD}, e' che il modello generativo profondo non debba apprendere
direttamente una distribuzione su matrici di adiacenza, ma una distribuzione
\emph{sui graphon}. Il vantaggio atteso e' duplice. In primo luogo, la
rappresentazione a graphon e' invariante rispetto alla dimensione del grafo: un
unico oggetto \(128\times 128\) puo' generare grafi di qualunque numero di nodi.
In secondo luogo, la rappresentazione e' gia' allineata (i nodi sono ordinati in
base alle variabili latenti stimate), il che elimina il problema della
permutazione dal modello generativo e lo riduce a un problema di modellazione di
un tensore denso a valori in \([0,1]\).

\section{Struttura della pipeline e contributi}

La pipeline si articola in quattro stadi, illustrati in dettaglio nel
Capitolo~\ref{ch:pipeline}:

\begin{enumerate}
    \item \textbf{Stima dei graphon.} A partire da un insieme di grafi osservati,
    SIGL (\emph{Scalable Implicit Graphon Learning}) stima graphon mediante
    rappresentazioni neurali implicite (INR) e reti neurali su grafi
    \cite{azizpour2025sigl}. Una modifica introdotta in questo lavoro consente di
    ottenere da un singolo dataset non un unico graphon ma un intero
    \emph{campione} di graphon, tramite clustering dei grafi seguito da bootstrap
    non parametrico.
    \item \textbf{Modellazione generativa dei graphon.} I graphon stimati sono
    usati come dataset di training per un modello di diffusione in spazio discreto.
    Sono implementate due varianti alternative e mutuamente esclusive:
    \begin{itemize}
        \item \textbf{OGAD--D3PM}, che tratta il graphon quantizzato come
        un'immagine categoriale \(128\times 128\) e applica il framework D3PM
        \cite{austin2021structured} con un denoiser di tipo Diffusion Transformer
        \cite{d3pm_pytorch};
        \item \textbf{OGAD--DiGress}, che tratta il graphon quantizzato come un
        \emph{grafo completo} su \(128\) nodi con etichette d'arco categoriali, e
        applica il modello DiGress \cite{vignac2023digress} senza modifiche al
        nucleo diffusivo.
    \end{itemize}
    \item \textbf{Campionamento dei grafi.} I graphon generati vengono sanificati
    e usati come kernel di Bernoulli per campionare grafi, con numero di nodi
    estratto dalla distribuzione empirica del test set.
    \item \textbf{Valutazione.} I grafi generati sono confrontati con il test set
    tramite MMD su descrittori strutturali, seguendo il protocollo di selezione di
    kernel e iperparametri proposto da GGME \cite{obray2022evaluation}.
\end{enumerate}

Il contributo principale di questa tesi non e' la proposta di un nuovo modello
diffusivo, ma la \emph{costruzione e la caratterizzazione precisa} di una pipeline
composita che collega quattro lavori indipendenti (SIGL, D3PM, DiGress, GGME) in
un protocollo sperimentale riproducibile, insieme all'analisi delle scelte di
progetto --- in particolare della discretizzazione dei valori del graphon --- che
determinano la qualita' del risultato finale.

\section{Obiettivi}

Gli obiettivi sperimentali possono essere riassunti come segue:

\begin{itemize}
    \item stimare un insieme di graphon rappresentativi del training set mediante
    SIGL, con una procedura di clustering e bootstrap che ne aumenti la
    cardinalita' fino a rendere possibile l'addestramento di un modello generativo;
    \item addestrare un generatore diffusivo discreto sui graphon stimati, nelle
    due varianti D3PM e DiGress;
    \item confrontare la pipeline OGAD con due riferimenti: DiGress applicato
    direttamente ai grafi (baseline), e i grafi indicati come \enquote{Paper},
    ossia le predizioni di esempio distribuite con il repository GGME;
    \item valutare i grafi generati tramite MMD su distribuzione dei gradi,
    coefficiente di clustering e spettro del laplaciano normalizzato.
\end{itemize}

\section{Fondamenti bibliografici}

Il lavoro si colloca all'intersezione tra teoria dei graphon, stima statistica di
reti e modelli generativi profondi per grafi. Dal punto di vista teorico, i
graphon nascono come oggetti limite per sequenze dense di grafi
\cite{lovasz2006limits,lovasz2012large}; la connessione con i modelli di grafo
scambiabili e con la rappresentazione di Aldous--Hoover e' discussa in modo
esplicito da Diaconis e Janson \cite{diaconis2008graph}. Questa base giustifica
l'uso del graphon come rappresentazione probabilistica non parametrica della
struttura di una famiglia di grafi.

Sul piano statistico, la stima di graphon e' stata studiata sia come problema
minimax sia come problema computazionale. Gao, Lu e Zhou stabiliscono tassi
ottimali per la stima di graphon e mettono in evidenza il legame con modelli a
blocchi stocastici e regressione non parametrica \cite{gao2015rate}. Borgs e
Chayes presentano invece una prospettiva piu' ampia sull'uso dei graphon per
modellare, stimare e progettare algoritmi su reti massive \cite{borgs2017graphons}.
Chan e Airoldi introducono lo stimatore a istogramma consistente su grafi ordinati
\cite{chan2014histogram} che costituisce il secondo stadio di SIGL, mentre IGNR
\cite{xia2023ignr} e SIGL \cite{azizpour2025sigl} sostituiscono la matrice a
risoluzione fissa con una rappresentazione neurale implicita continua.

Per quanto riguarda la generazione di grafi, GraphRNN rappresenta un riferimento
classico per i modelli autoregressivi e ha contribuito a rendere comune la
valutazione tramite MMD su descrittori strutturali \cite{you2018graphrnn}.
Altri lavori, come GraphVAE e MolGAN, mostrano approcci alternativi basati
rispettivamente su autoencoder variazionali e modelli avversari
\cite{simonovsky2018graphvae,de2018molgan}. La componente diffusiva della tesi
si collega ai diffusion probabilistic models \cite{sohl2015deep,ho2020denoising},
alla formulazione score-based tramite equazioni differenziali stocastiche
\cite{song2021scorebased}, e alla loro estensione a spazi discreti tramite
multinomial diffusion \cite{hoogeboom2021argmax} e D3PM
\cite{austin2021structured}, di cui DiGress \cite{vignac2023digress} costituisce
la specializzazione al dominio dei grafi.

\chapter{Contesto teorico}
\label{ch:teoria}

Questo capitolo raccoglie la formalizzazione dei quattro componenti della
pipeline. L'obiettivo e' fissare la notazione che verra' usata nel
Capitolo~\ref{ch:pipeline}, dove ciascun componente viene specializzato al caso
concreto implementato nel repository.

\section{Graphon e generazione di grafi}
\label{sec:graphon}

\begin{definition}[Graphon]
Un graphon e' una funzione misurabile, simmetrica e limitata
\[
    W : [0,1]^2 \to [0,1],
    \qquad
    W(u,v) = W(v,u)
    \quad\text{per quasi ogni } (u,v).
\]
\end{definition}

Dato un intero \(M \in \mathbb{N}\), il graphon induce una distribuzione su grafi
semplici non orientati con \(M\) nodi. Il modello generativo associato, detto
\emph{\(W\)-random graph}, e'
\begin{equation}
    \label{eq:graphon-sampling}
    \begin{aligned}
        \eta_1,\ldots,\eta_M
        &\overset{\mathrm{i.i.d.}}{\sim} \Uniform([0,1]),\\
        A_{ij} \mid \eta_i,\eta_j,W
        &\sim \Bernoulli\bigl(W(\eta_i,\eta_j)\bigr),
        \qquad 1 \leq i < j \leq M,
    \end{aligned}
\end{equation}
con \(A_{ji}=A_{ij}\) e \(A_{ii}=0\). Condizionatamente alle variabili latenti dei
nodi, gli archi sono indipendenti e la probabilita' di un arco dipende solo dal
valore del graphon nella coppia di coordinate latenti:
\[
    \GraphDist(A_{ij}=1 \mid \eta_i,\eta_j,W) = W(\eta_i,\eta_j).
\]
Il teorema di Aldous--Hoover garantisce che ogni distribuzione scambiabile su
grafi infiniti ammette una rappresentazione di questa forma \cite{diaconis2008graph},
il che giustifica l'uso del graphon come modello non parametrico universale per
famiglie di grafi scambiabili.

\subsection{Non identificabilita' e problema dell'ordinamento}

La rappresentazione~\eqref{eq:graphon-sampling} non e' identificabile. Se
\(\varphi:[0,1]\to[0,1]\) e' una trasformazione misura-preservante, allora
\begin{equation}
    \label{eq:graphon-equivalence}
    W'(u,v) = W(\varphi(u),\varphi(v))
    \quad \text{q.o.}
\end{equation}
induce esattamente la stessa distribuzione su grafi. Ne segue che la classe di
equivalenza \([W]\), e non \(W\) stesso, e' l'oggetto identificabile.

Questo fatto ha due conseguenze operative centrali per la pipeline OGAD:

\begin{enumerate}
    \item \textbf{Lato stima.} Uno stimatore di graphon deve fissare un
    rappresentante della classe, cioe' scegliere un \emph{ordinamento} dei nodi.
    I metodi classici (SAS, USVT) ordinano per grado empirico
    \cite{chan2014histogram,azizpour2025sigl}; SIGL apprende invece l'ordinamento
    tramite una GNN, il che gli consente di trattare anche graphon non monotoni
    come i modelli a blocchi stocastici, dove tutti i nodi hanno lo stesso grado
    atteso.
    \item \textbf{Lato generazione.} Poiche' il modello diffusivo apprende una
    distribuzione su \emph{matrici} \(128\times128\) e non su classi di
    equivalenza, la coerenza dell'ordinamento prodotto da SIGL sui diversi graphon
    del training set e' un prerequisito: se due graphon equivalenti fossero
    rappresentati con permutazioni latenti diverse, il modello dovrebbe spendere
    capacita' per apprendere l'orbita della permutazione anziche' la struttura.
\end{enumerate}

\subsection{Discretizzazione}

Nell'implementazione, \(W\) non e' rappresentato come funzione continua ma come
matrice \(R \times R\), con \(R=128\). Indicando con
\(\widehat{W} \in [0,1]^{R \times R}\) tale discretizzazione, il campionamento
di~\eqref{eq:graphon-sampling} viene approssimato scegliendo indici
\begin{equation}
    \label{eq:latent-index}
    z_i = \bigl\lfloor R\,\eta_i \bigr\rfloor \wedge (R-1),
    \qquad z_i \in \{0,\ldots,R-1\},
\end{equation}
e ponendo
\[
    A_{ij} \mid z_i,z_j,\widehat{W}
    \sim \Bernoulli\bigl(\widehat{W}_{z_i z_j}\bigr).
\]
Prima del campionamento si applica una proiezione
\(\Pi : \Real^{R \times R} \to [0,1]^{R \times R}\) che impone i vincoli
strutturali di grafo semplice non orientato:
\begin{equation}
    \label{eq:sanitize}
    \Pi(X)_{ij}
    =
    \begin{cases}
        0, & i=j,\\[2pt]
        \min\bigl\{1,\max\{0,(X_{ij}+X_{ji})/2\}\bigr\}, & i \neq j.
    \end{cases}
\end{equation}
Ossia: simmetrizzazione, clipping in \([0,1]\), diagonale nulla. Questa proiezione
ricorre in tre punti distinti della pipeline (uscita di SIGL, uscita del
generatore, ingresso del campionatore) ed e' l'unico meccanismo che garantisce la
validita' strutturale del graphon, dal momento che nessuno dei due modelli
generativi impone la simmetria a livello architetturale.

\section{Stima del graphon: SIGL}
\label{sec:sigl-theory}

SIGL \cite{azizpour2025sigl} risolve il seguente problema: dato un insieme di
grafi \(\mathcal{D}=\{A_t\}_{t=1}^{M}\) generati da un graphon incognito \(\omega\),
stimare \(\omega\). Rispetto a IGNR \cite{xia2023ignr}, che minimizza la distanza
di Gromov--Wasserstein e ha complessita' cubica nel numero di nodi, SIGL
sostituisce la GW con un MSE puntuale, il che richiede pero' di risolvere
esplicitamente il problema di allineamento. La procedura si articola in tre passi.

\paragraph{Passo 1: apprendimento dell'ordinamento.}
Si stimano le variabili latenti \(\hat{\eta}_t = \{\hat\eta_i\}_{i=1}^{n_t}\) di
ogni grafo mediante una GNN \(g_{\phi_1}\), e simultaneamente un
\emph{graphon ausiliario} \(h_{\phi_2}\), rappresentato da una INR, che funge da
riferimento comune per tutti i grafi:
\begin{equation}
    \label{eq:sigl-step1}
    \mathcal{L}(\phi)
    =
    \sum_{t=1}^{M}
    \frac{1}{n_t^2}
    \sum_{i,j=1}^{n_t}
    \bigl[
        A_t(i,j) - h_{\phi_2}\bigl(\hat\eta_t(i),\hat\eta_t(j)\bigr)
    \bigr]^2 ,
    \qquad
    \hat{\eta}_t = g_{\phi_1}(A_t, X_t).
\end{equation}
Le feature di nodo \(X_t\) sono rumore gaussiano, cosicche' l'informazione usata
per l'ordinamento e' interamente strutturale. Si definisce quindi la permutazione
\(\hat\pi\) che ordina i nodi per valore decrescente di \(\hat\eta\), e si ottiene
la matrice di adiacenza ordinata \(\hat{A}_t(i,j) = A_t(\hat\pi(i),\hat\pi(j))\).

\paragraph{Passo 2: approssimazione a istogramma.}
Seguendo Chan e Airoldi \cite{chan2014histogram}, si applica un average pooling non
sovrapposto di finestra \(h\) alla matrice ordinata:
\begin{equation}
    \label{eq:sigl-hist}
    \hat{H}_t(i,j)
    =
    \frac{1}{h^2}
    \sum_{s_1=1}^{h}\sum_{s_2=1}^{h}
    \hat{A}_t\bigl((i-1)h+s_1,\;(j-1)h+s_2\bigr),
    \qquad
    \hat{H}_t \in [0,1]^{k_t \times k_t},
\end{equation}
con \(k_t = \lfloor n_t / h \rfloor\). Ciascun \(\hat{H}_t\) e' una realizzazione
rumorosa del graphon vero, osservata a una risoluzione diversa.

\paragraph{Passo 3: regressione INR.}
Si costruisce il dataset supervisionato
\[
    \mathcal{C}
    =
    \Bigl\{
        \bigl(\tfrac{i}{k_t},\tfrac{j}{k_t},\hat{H}_t(i,j)\bigr)
        \;:\;
        i,j \in \{1,\ldots,k_t\},\;
        t \in \{1,\ldots,M\}
    \Bigr\},
\]
e si addestra una INR \(f_\theta : [0,1]^2 \to [0,1]\) minimizzando un MSE pesato:
\begin{equation}
    \label{eq:sigl-step3}
    \theta^\star
    =
    \argmin_\theta
    \sum_{t=1}^{M} w_t
    \sum_{i=1}^{k_t}\sum_{j=1}^{k_t}
    \left(
        \hat{H}_t(i,j)
        - f_\theta\!\left(\tfrac{i}{k_t},\tfrac{j}{k_t}\right)
    \right)^2 ,
    \qquad
    w_t = \frac{|G_t|}{\sum_{i=1}^{M}|G_i|}.
\end{equation}
Il peso \(w_t\) proporzionale al numero di nodi privilegia i grafi piu' grandi, la
cui risoluzione fornisce una stima migliore del graphon vero. La INR e' una rete
SIREN \cite{sitzmann2020siren}, cioe' un MLP con attivazioni sinusoidali
\(\phi_i(x) = \sin(\omega_0 W_i x + b_i)\), scelta per la capacita' di
rappresentare segnali ad alta frequenza. Essendo \(f_\theta\) definita
sull'intero dominio continuo, il graphon stimato e' \emph{resolution-free}: puo'
essere campionato a qualunque risoluzione \(R\).

La Proposizione~1 di SIGL stabilisce che, se \(\omega\) e' Lipschitz con costante
\(\mathcal{L}\) e si sceglie \(h = \log n\), allora
\begin{equation}
    \label{eq:sigl-consistency}
    \lim_{n\to\infty}
    \mathbb{E}\bigl[\|f_\theta - \omega\|_2^2\bigr]
    \leq
    \varepsilon_{\mathrm{tr}} + 4\mathcal{L}^2\tau^2,
\end{equation}
dove \(\varepsilon_{\mathrm{tr}}\) e' l'errore massimo di regressione della INR sui
punti di \(\mathcal{C}\) e \(\tau\) e' l'errore massimo della GNN nella stima delle
variabili latenti. L'errore asintotico e' quindi controllato \emph{solo} dalla
capacita' espressiva dei due moduli: architetture piu' espressive producono
stimatori piu' consistenti.

\section{Diffusione in spazi discreti: il framework D3PM}
\label{sec:d3pm-theory}

I modelli di diffusione apprendono a invertire un processo di corruzione
progressiva dei dati \cite{sohl2015deep,ho2020denoising}. Nei modelli continui la
corruzione e' aggiunta di rumore gaussiano; in spazi discreti occorre invece
definire transizioni tra stati categoriali. D3PM \cite{austin2021structured}
fornisce il framework generale.

\subsection{Processo diretto}

Sia \(x_0 \in \{1,\ldots,K\}\) una variabile categoriale e \(\boldsymbol{x}_0\) la
sua codifica one-hot come vettore riga. Il processo diretto e' una catena di Markov
\begin{equation}
    q(x_{1:T}\mid x_0) = \prod_{t=1}^{T} q(x_t \mid x_{t-1}),
    \qquad
    q(x_t \mid x_{t-1}) = \Cat\bigl(x_t;\, \boldsymbol{x}_{t-1} Q_t\bigr),
\end{equation}
dove \([Q_t]_{ij} = q(x_t=j \mid x_{t-1}=i)\) e' una matrice stocastica per riga.
Poiche' la catena e' markoviana, la marginale a \(t\) passi e la posterior sono
disponibili in forma chiusa:
\begin{align}
    \label{eq:d3pm-marginal}
    q(x_t\mid x_0) &= \Cat\bigl(x_t;\, \boldsymbol{x}_0 \overline{Q}_t\bigr),
    \qquad \overline{Q}_t = Q_1 Q_2 \cdots Q_t,\\[4pt]
    \label{eq:d3pm-posterior}
    q(x_{t-1}\mid x_t, x_0)
    &=
    \Cat\!\left(
        x_{t-1};\,
        \frac{\boldsymbol{x}_t Q_t^{\top} \odot \boldsymbol{x}_0 \overline{Q}_{t-1}}
             {\boldsymbol{x}_0 \overline{Q}_t \boldsymbol{x}_t^{\top}}
    \right),
\end{align}
con \(\odot\) prodotto puntuale. Sono le due proprieta' che rendono possibile
l'addestramento efficiente: \eqref{eq:d3pm-marginal} consente di campionare
direttamente \(x_t\) da \(x_0\) per un \(t\) arbitrario (training parallelo sui
timestep), \eqref{eq:d3pm-posterior} consente di calcolare analiticamente le
divergenze KL della loss.

\subsection{Scelta della matrice di transizione}

D3PM mostra che la scelta di \(Q_t\) e' una decisione di progetto rilevante.
Le due varianti pertinenti a questo lavoro sono:

\begin{itemize}
    \item \textbf{Uniforme} \cite{hoogeboom2021argmax}:
    \begin{equation}
        \label{eq:q-uniform}
        Q_t = (1-\beta_t) I + \frac{\beta_t}{K}\,\one\one^{\top},
        \qquad
        [Q_t]_{ij} =
        \begin{cases}
            1 - \frac{K-1}{K}\beta_t, & i=j,\\[2pt]
            \frac{\beta_t}{K}, & i \neq j.
        \end{cases}
    \end{equation}
    La matrice e' doppiamente stocastica a entrate strettamente positive, quindi
    per il teorema di Perron--Frobenius la distribuzione stazionaria e' uniforme
    su \(\{1,\ldots,K\}\). E' questa la variante usata in \emph{entrambe} le
    configurazioni OGAD.
    \item \textbf{Marginale} \cite{vignac2023digress}:
    \(Q_t = \alpha^t I + \beta^t \one \boldsymbol{m}^{\top}\), con
    \(\boldsymbol{m}\) la distribuzione marginale empirica delle classi nel
    training set. La distribuzione limite e' \(\boldsymbol{m}\) anziche' uniforme;
    l'idea e' che una prior piu' vicina alla distribuzione dei dati renda il
    training piu' facile. Il Teorema 4.1 di DiGress mostra che il prodotto delle
    marginali e' la proiezione ortogonale della distribuzione vera sulla classe
    delle distribuzioni prodotto.
\end{itemize}

\subsection{Schedule di rumore}

Lo schedule determina \(\beta_t\). La pipeline usa in entrambe le varianti uno
schedule di tipo coseno, ispirato a Nichol e Dhariwal \cite{nichol2021improved},
che fissa la probabilita' cumulata di non-transizione \(\bar\alpha_t\) e ne deduce
\(\beta_t\):
\begin{equation}
    \label{eq:cosine}
    \bar\alpha(s) = \cos\!\left(\frac{s+s_0}{1+s_0}\cdot\frac{\pi}{2}\right)^{p},
    \qquad
    \beta_t = 1 - \frac{\bar\alpha(t/T)}{\bar\alpha((t-1)/T)},
    \qquad s_0 = 0.008 .
\end{equation}
Le due implementazioni differiscono nell'esponente \(p\) (Sezione~\ref{sec:d3pm-impl}
e Sezione~\ref{sec:digress-impl}): D3PM nel repository usa \(p=1\), DiGress usa
\(p=2\) come nella formulazione originale.

\subsection{Parametrizzazione inversa e loss}

Invece di predire direttamente i logit di \(p_\theta(x_{t-1}\mid x_t)\), si segue
la \(x_0\)-parametrizzazione: una rete \(\mathrm{nn}_\theta\) predice i logit di
\(\tilde{p}_\theta(\tilde{x}_0 \mid x_t)\), che vengono combinati con la posterior
esatta:
\begin{equation}
    \label{eq:x0-param}
    p_\theta(x_{t-1}\mid x_t)
    \;\propto\;
    \sum_{\tilde{x}_0}
    q(x_{t-1}, x_t \mid \tilde{x}_0)\,
    \tilde{p}_\theta(\tilde{x}_0\mid x_t).
\end{equation}
Questa scelta ha tre vantaggi: (i) la KL in \eqref{eq:d3pm-elbo} si annulla
esattamente quando \(\tilde{p}_\theta\) concentra tutta la massa su \(x_0\);
(ii) il pattern di sparsita' di \(Q_t\) e' automaticamente ereditato da
\(p_\theta\); (iii) consente inferenza a \(k\) passi alla volta,
\(p_\theta(x_{t-k}\mid x_t) = \sum q(x_{t-k},x_t\mid \tilde{x}_0)
\tilde{p}_\theta(\tilde{x}_0\mid x_t)\).

L'addestramento minimizza un bound variazionale sulla negative log-likelihood:
\begin{equation}
    \label{eq:d3pm-elbo}
    L_{\mathrm{vb}}
    =
    \mathbb{E}_{q(x_0)}
    \Biggl[
        \underbrace{\KL\bigl(q(x_T\mid x_0)\,\|\,p(x_T)\bigr)}_{L_T}
        +
        \sum_{t=2}^{T}
        \underbrace{\mathbb{E}_{q(x_t\mid x_0)}
        \KL\bigl(q(x_{t-1}\mid x_t,x_0)\,\|\,p_\theta(x_{t-1}\mid x_t)\bigr)}_{L_{t-1}}
        \underbrace{-\,\mathbb{E}_{q(x_1\mid x_0)}\log p_\theta(x_0\mid x_1)}_{L_0}
    \Biggr].
\end{equation}
Il termine \(L_T\) non dipende da \(\theta\) se la prior e' fissata alla
distribuzione stazionaria del processo diretto. D3PM introduce inoltre una loss
ibrida che aggiunge un termine ausiliario di cross-entropy sulla predizione di
\(x_0\):
\begin{equation}
    \label{eq:d3pm-hybrid}
    L_\lambda
    =
    L_{\mathrm{vb}}
    + \lambda\,
    \mathbb{E}_{q(x_0)}\mathbb{E}_{q(x_t\mid x_0)}
    \bigl[-\log \tilde{p}_\theta(x_0\mid x_t)\bigr].
\end{equation}
Il termine ausiliario coincide con \(L_0\) per \(t=1\) ed e' minimizzato dalla
stessa configurazione che minimizza \(L_{\mathrm{vb}}\); la sua funzione e'
stabilizzare il gradiente, la cui varianza in \(L_{\mathrm{vb}}\) dipende
fortemente da \(t\). Nel repository \(\lambda\) e' l'iperparametro
\texttt{hybrid\_loss\_coeff}.

Nel caso multivariato --- che e' il caso di interesse, dato che un graphon ha
\(R^2\) celle --- si assume che \(Q_t\) sia applicata indipendentemente a ogni
componente e che sia \(q\) sia \(p_\theta\) fattorizzino sulle componenti. Le KL
si calcolano dunque componente per componente e si sommano.

\section{Diffusione discreta su grafi: DiGress}
\label{sec:digress-theory}

DiGress \cite{vignac2023digress} specializza il framework precedente al caso in
cui il dato e' un grafo con attributi categoriali su nodi e archi. Un grafo con
\(n\) nodi e' rappresentato dalla coppia \(G = (X, E)\), con
\(X \in \Real^{n\times a}\) matrice delle codifiche one-hot dei tipi di nodo
(\(a\) classi) ed \(E \in \Real^{n\times n\times b}\) tensore delle codifiche
one-hot dei tipi di arco (\(b\) classi, dove l'assenza di arco e' trattata come un
tipo particolare).

\subsection{Processo diretto fattorizzato}

Il rumore viene applicato indipendentemente a ogni nodo e a ogni arco, tramite due
matrici di transizione distinte \(Q_X^t \in \Real^{a\times a}\) e
\(Q_E^t \in \Real^{b\times b}\):
\begin{equation}
    \label{eq:digress-forward}
    q(G^t \mid G^{t-1}) = \bigl(X^{t-1}Q_X^t,\; E^{t-1}Q_E^t\bigr),
    \qquad
    q(G^t \mid G) = \bigl(X\overline{Q}_X^t,\; E\overline{Q}_E^t\bigr).
\end{equation}
Nel caso di grafi non orientati il rumore e' applicato solo alla parte
triangolare superiore di \(E\), che viene poi simmetrizzata. Si noti che lo
spazio degli stati su cui e' definita la matrice di transizione non e' lo spazio
dei grafi (che sarebbe di cardinalita' proibitiva) ma solo l'insieme dei tipi di
nodo e di arco.

\subsection{Rete di denoising e loss}

La rete \(\varphi_\theta\) riceve \(G^t\) e predice una distribuzione
\(\hat{p}^G = (\hat{p}^X, \hat{p}^E)\) sul grafo pulito. L'addestramento
minimizza una cross-entropy che si decompone su nodi e archi:
\begin{equation}
    \label{eq:digress-loss}
    l(\hat{p}^G, G)
    =
    \sum_{1\leq i\leq n} \mathrm{CE}\bigl(x_i, \hat{p}^X_i\bigr)
    \;+\;
    \lambda \sum_{1\leq i,j\leq n} \mathrm{CE}\bigl(e_{ij}, \hat{p}^E_{ij}\bigr),
\end{equation}
con \(\lambda \in \Real^{+}\) che bilancia il peso relativo di nodi e archi.
Il campionamento inverso marginalizza sulle predizioni della rete:
\begin{equation}
    \label{eq:digress-reverse}
    p_\theta(G^{t-1}\mid G^t)
    =
    \prod_{i} p_\theta(x_i^{t-1}\mid G^t)
    \prod_{i,j} p_\theta(e_{ij}^{t-1}\mid G^t),
\end{equation}
\begin{equation}
    p_\theta(x_i^{t-1}\mid G^t)
    =
    \sum_{x\in\mathcal{X}}
    q\bigl(x_i^{t-1}\mid x_i = x,\, x_i^t\bigr)\,\hat{p}^X_i(x),
\end{equation}
e analogamente per gli archi.

\subsection{Proprieta' di equivarianza}

Tre lemmi di DiGress garantiscono la correttezza rispetto alle permutazioni:
l'architettura e' equivariante (Lemma 3.1), la loss~\eqref{eq:digress-loss} e'
invariante perche' decomponibile per nodo e per arco (Lemma 3.2), e la
distribuzione generata e' scambiabile (Lemma 3.3), poiche' la prior e' un prodotto
di distribuzioni i.i.d. e le transizioni sono equivarianti. In particolare il
Lemma 3.2 implica che non e' necessario alcun graph matching tra grafo predetto e
grafo target, a differenza degli approcci basati su VAE.

\begin{remark}
Nella variante OGAD--DiGress (Sezione~\ref{sec:digress-impl}) queste proprieta'
restano valide ma perdono rilevanza pratica: il \enquote{grafo} in ingresso e' un
grafo completo su \(128\) nodi con tipi di nodo costanti, quindi l'unica
informazione risiede nelle etichette d'arco e la permutazione dei nodi corrisponde
proprio alla riparametrizzazione~\eqref{eq:graphon-equivalence} che SIGL ha gia'
fissato. L'equivarianza della rete diventa quindi un vincolo che il modello deve
\enquote{superare}: la struttura del graphon e' \emph{ordinata}, e un modello
equivariante non puo' sfruttare direttamente tale ordinamento se non attraverso
i valori delle etichette.
\end{remark}

\subsection{Feature ausiliarie strutturali}

Poiche' le MPNN standard hanno potere espressivo limitato (non sanno contare i
cicli), DiGress aumenta l'input della rete con descrittori calcolati sul grafo
rumoroso: conteggi di cicli fino a lunghezza 6, numero di componenti connesse,
primi autovalori del laplaciano e relativi autovettori. Questa possibilita' e'
resa disponibile proprio dalla natura discreta del rumore, che preserva la
sparsita' del grafo rumoroso (a differenza della diffusione gaussiana, che
produce grafi completi in cui le nozioni topologiche non sono definite).

\section{Valutazione tramite MMD e il protocollo GGME}
\label{sec:ggme-theory}

La valutazione dei modelli generativi per grafi e' particolarmente delicata:
diversamente dalle immagini, i grafi non possono essere valutati per ispezione
visiva, e le metriche comunemente adottate risultano instabili se non associate a
descrittori strutturali e kernel appropriati. GGME \cite{obray2022evaluation}
affronta sistematicamente questo problema.

\subsection{Desiderata}

GGME enuncia tre requisiti per una metrica di confronto tra distribuzioni di grafi:
\begin{itemize}
    \item \textbf{Espressivita'}: \(d(\Graphs,\Graphs')\) deve crescere
    monotonicamente al crescere della dissimilarita';
    \item \textbf{Robustezza}: piccole perturbazioni della distribuzione devono
    produrre piccole variazioni della metrica;
    \item \textbf{Efficienza}: il costo computazionale deve scalare con il numero e
    la dimensione dei grafi.
\end{itemize}

\subsection{Descrittori e MMD}

Siano \(\mathcal{P}\) e \(\mathcal{Q}\) due distribuzioni su grafi e sia
\(\phi : \Graphs \to \Real^m\) un descrittore strutturale. Data una funzione
kernel positiva definita \(k : \Real^m \times \Real^m \to \Real\), la MMD tra le
distribuzioni indotte dai descrittori e'
\begin{equation}
    \label{eq:mmd}
    \begin{aligned}
        \MMD^2_k(\mathcal{P},\mathcal{Q})
        &=
        \mathbb{E}_{G,G'\sim\mathcal{P}}\,k(\phi(G),\phi(G'))
        +
        \mathbb{E}_{H,H'\sim\mathcal{Q}}\,k(\phi(H),\phi(H'))\\
        &\quad-
        2\,\mathbb{E}_{G\sim\mathcal{P},H\sim\mathcal{Q}}\,k(\phi(G),\phi(H)).
    \end{aligned}
\end{equation}
Dato un campione osservato \(\{G_i\}_{i=1}^{n}\) e uno generato
\(\{H_j\}_{j=1}^{m}\), lo stimatore \emph{non distorto}
\cite{gretton2012kernel} rimuove i termini diagonali di auto-confronto:
\begin{equation}
    \label{eq:mmd-unbiased}
    \begin{aligned}
        \widehat{\MMD}^2_k
        &=
        \frac{1}{n(n-1)}
        \sum_{i \neq i'} k(\phi(G_i),\phi(G_{i'}))
        +
        \frac{1}{m(m-1)}
        \sum_{j \neq j'} k(\phi(H_j),\phi(H_{j'}))\\
        &\quad-
        \frac{2}{nm}
        \sum_{i=1}^{n}\sum_{j=1}^{m}
            k(\phi(G_i),\phi(H_j)).
    \end{aligned}
\end{equation}
Essendo non distorto, lo stimatore \eqref{eq:mmd-unbiased} puo' assumere valori
negativi; e' inoltre uno stimatore di \(\MMD^2\), non di \(\MMD\).

I kernel considerati nella pipeline sono il kernel lineare
\(k(a,b) = \langle a,b\rangle\), privo di iperparametri, e il kernel gaussiano
\begin{equation}
    \label{eq:rbf}
    k_\sigma(a,b) = \exp\bigl(-\sigma\,\|a-b\|_2^2\bigr).
\end{equation}
Si noti che qui \(\sigma\) parametrizza una \emph{precisione} e non una larghezza
di banda. GGME sconsiglia esplicitamente il kernel basato su Earth Mover's Distance
(troppo costoso, fino a due ordini di grandezza piu' lento) e il kernel gaussiano
basato sulla distanza di variazione totale, che non e' definito positivo perche' lo
spazio metrico indotto dalla TV non e' piatto (Teorema 1 di GGME).

\subsection{Selezione di kernel e iperparametri}

Il contributo metodologico principale di GGME e' la dimostrazione che il ranking
dei modelli e' \emph{arbitrario} rispetto alla scelta di kernel, \(\sigma\) e
numero di bin del descrittore: qualunque modello puo' risultare primo con una
scelta opportuna. La soluzione proposta e' selezionare tali elementi mediante un
esperimento di perturbazione controllata che \emph{non coinvolge i modelli}:

\begin{enumerate}
    \item si genera una sequenza di versioni perturbate del test set con intensita'
    crescente \(p_1 < p_2 < \cdots < p_L\);
    \item per ogni combinazione (kernel, \(\sigma\)) si calcola
    \(\widehat{\MMD}^2\) tra test set e ciascuna versione perturbata;
    \item si sceglie la combinazione che massimizza la correlazione di Pearson tra
    l'intensita' della perturbazione e la MMD risultante.
\end{enumerate}

La procedura e' non distorta rispetto ai modelli in confronto, perche' dipende
solo dalla distribuzione di partenza. GGME raccomanda inoltre di riportare la MMD
tra training e test set come \emph{scala di riferimento}, dato che la MMD non ha
una scala intrinseca.

\chapter{La pipeline OGAD}
\label{ch:pipeline}

Questo capitolo descrive la pipeline implementata, collegando ciascuna scelta
implementativa al corrispondente formalismo del Capitolo~\ref{ch:teoria}. La
struttura complessiva e'
\[
    \underbrace{\text{SIGL}}_{\text{stadio 1}}
    \;\longrightarrow\;
    \underbrace{\text{D3PM \emph{oppure} DiGress su graphon}}_{\text{stadio 2}}
    \;\longrightarrow\;
    \underbrace{\text{campionamento di grafi}}_{\text{stadio 3}}
    \;\longrightarrow\;
    \underbrace{\text{GGME}}_{\text{stadio 4}} .
\]

Formalmente, indicando con
\(\mathcal{D}_{\mathrm{train}} = \{G_1,\ldots,G_N\}\) il training set di grafi, il
primo stadio produce una collezione di graphon stimati
\begin{equation}
    \widehat{\mathcal{W}}
    =
    S_{\psi}(\mathcal{D}_{\mathrm{train}})
    =
    \{\widehat{W}_1,\ldots,\widehat{W}_L\},
    \qquad \widehat{W}_\ell \in [0,1]^{R\times R},
\end{equation}
dove \(S_\psi\) rappresenta lo stimatore SIGL. Il secondo stadio apprende una
distribuzione \(p_\theta(W)\) su tale collezione e ne estrae campioni
\(\widetilde{W}_1,\ldots,\widetilde{W}_S \sim p_\theta\). Il terzo stadio induce
una distribuzione su grafi
\begin{equation}
    \widetilde{G} \sim \GraphDist\bigl(\,\cdot \mid \Pi(\widetilde{W}), M\bigr),
    \qquad M \sim \hat{\mu}_M,
\end{equation}
con \(\hat{\mu}_M\) la distribuzione empirica delle dimensioni dei grafi di test.
Il quarto stadio confronta
\(\GraphDist_{\Graphs}^{\mathrm{test}}\) con
\(\GraphDist_{\Graphs,\theta,\psi}^{\mathrm{OGAD}}\).

\section{Dataset}
\label{sec:dataset}

Tutti gli esperimenti usano il dataset \emph{Community Graphs}, distribuito con il
repository GGME e memorizzato in formato pickle:

\begin{center}
\begin{tabular}{lrrl}
    \toprule
    File & \(N\) grafi & Nodi & Formato \\
    \midrule
    \texttt{data/train/graphs/CommunityGraphs.pickle} & 400 & 45 & \texttt{csr\_matrix} \\
    \texttt{data/test/graphs/CommunityGraphs.pickle} & 100 & 45 & \texttt{csr\_matrix} \\
    \texttt{data/predictions/graphs/ggme\_paper/} & 400 & 44--45 & \texttt{ndarray} \\
    \bottomrule
\end{tabular}
\end{center}

La densita' media di archi e' \(\approx 0.234\) sul training set e
\(\approx 0.232\) sul test set. Il fatto che \emph{tutti} i grafi abbiano
esattamente 45 nodi ha una conseguenza importante sul terzo stadio: la
distribuzione empirica \(\hat\mu_M\) e' degenere in \(45\), quindi tutti i grafi
generati hanno 45 nodi e le differenze nelle metriche non possono essere imputate
a una diversa distribuzione delle dimensioni.

Il terzo file, indicato nel confronto come \enquote{Paper}, contiene le predizioni
di esempio distribuite con il repository GGME: sono i grafi prodotti da un modello
generativo di riferimento, e fungono da soglia di qualita' esterna.

\section{Stadio 1: stima dei graphon con SIGL}
\label{sec:sigl-impl}

\subsection{Il problema della cardinalita' del training set}

SIGL, come formulato in \cite{azizpour2025sigl}, produce \emph{un solo} graphon a
partire da un insieme di grafi. Questo e' incompatibile con lo stadio 2, che
richiede un dataset di graphon sufficientemente numeroso per addestrare un modello
di diffusione. La pipeline OGAD risolve il problema con una procedura in due
passi: \emph{clustering} dei grafi seguito da \emph{bootstrap non parametrico}
all'interno di ciascun cluster.

\paragraph{Clustering.}
Ogni grafo \(G_i\) e' mappato in un vettore di feature ottenuto concatenando i
\emph{medesimi tre descrittori} usati nella valutazione GGME, in forma
normalizzata:
\begin{equation}
    \label{eq:cluster-features}
    \mathbf{f}(G) =
    \bigl[\,
        \phi_{\mathrm{deg}}(G)
        \;\Vert\;
        \phi_{\mathrm{clust}}(G)
        \;\Vert\;
        \phi_{\mathrm{lap}}(G)
    \,\bigr].
\end{equation}
I vettori vengono zero-padded a lunghezza comune, standardizzati, e partizionati
con \(k\)-means in \(C\) cluster (\(C=50\) nella configurazione principale,
\(C=200\) in un test preliminare). Si ottiene una partizione
\(\{\mathcal{I}_1,\ldots,\mathcal{I}_C\}\) degli indici dei grafi.

\paragraph{Bootstrap.}
Per ogni cluster \(c\) e per ogni replica \(r = 1,\ldots,B\) si estrae un campione
bootstrap di indici
\begin{equation}
    \label{eq:bootstrap}
    \mathcal{I}_c^{(r)}
    =
    \bigl(
        i_1^{(r)},\ldots,i_{|\mathcal{I}_c|}^{(r)}
    \bigr),
    \qquad
    i_s^{(r)} \overset{\mathrm{i.i.d.}}{\sim} \Uniform(\mathcal{I}_c),
\end{equation}
cioe' \(|\mathcal{I}_c|\) estrazioni \emph{con reinserimento} dai membri del
cluster. Il numero totale di graphon prodotti e'
\begin{equation}
    \label{eq:L-count}
    L =
    \begin{cases}
        C, & B = 1,\\
        C \cdot B + C\,\mathbb{1}[\text{include\_original}], & B > 1.
    \end{cases}
\end{equation}
Con \(C=50\) e \(B \in \{12,13,14,15,16\}\) si ottengono \(L \in \{600, 650, 700,
750, 800\}\) graphon. Ogni job di bootstrap innesca una stima SIGL
\emph{completamente indipendente}: sia la GCN sia la SIREN sono reinizializzate da
zero, cosicche' i \(L\) graphon sono \(L\) stime distinte e non \(L\) viste dello
stesso oggetto.

\begin{remark}[Interpretazione statistica]
La procedura puo' essere letta come una stima della \emph{distribuzione
campionaria} dello stimatore SIGL, condizionata alla struttura locale del
dataset. Il clustering identifica sottopopolazioni omogenee rispetto ai descrittori
strutturali; il bootstrap all'interno di ciascuna genera la variabilita' che il
modello diffusivo deve poi apprendere. Senza bootstrap, il modello vedrebbe \(C\)
graphon quasi deterministici; senza clustering, la variabilita' bootstrap sarebbe
dominata dall'eterogeneita' globale del dataset.
\end{remark}

\paragraph{Split di validazione.}
Prima del clustering, il training set viene suddiviso in modo deterministico
(seed \(0\), frazione \(0.2\)) in \(320\) grafi usati per la stima dei graphon e
\(80\) grafi tenuti da parte come \emph{validation set GGME}, impiegato nello
stadio 4 per la selezione del candidato migliore (Sezione~\ref{sec:selection}).
Questo evita che la selezione del modello usi il test set.

\subsection{Iperparametri e architettura effettivi}

\begin{center}
\begin{tabular}{lll}
    \toprule
    Componente & Iperparametro & Valore \\
    \midrule
    GCN \(g_{\phi_1}\) & strati \texttt{GCNConv} & 3, con 8 unita' ciascuno \\
      & attivazione & ReLU, poi \(\mathrm{Linear}(8\to1)\) + sigmoide \\
      & normalizzazione & min--max per grafo, in \([0,1]\) \\
      & feature di nodo & \(x_i \sim \mathcal{N}(0,1)\), scalari \\
    SIREN \(h_{\phi_2}, f_\theta\) & strati nascosti & 2, con 20 unita' \\
      & \(\omega_0\) primo strato & 10 \\
      & \(\omega_0\) strati successivi & 1 \\
      & attivazione finale & sigmoide \\
    Ottimizzazione & epoche (passi 1 e 3) & 100 \\
      & learning rate & 0.06 (Adam) \\
      & batch (passo 1) & 1 grafo \\
      & batch (passo 3) & 1024 coordinate \\
    Discretizzazione & finestra di pooling \(h\) & \(\lfloor \ln n \rfloor = 3\) \\
      & risoluzione istogramma \(k\) & \(\lfloor 45/3\rfloor = 15\) \\
      & risoluzione di uscita \(R\) & 128 \\
    \bottomrule
\end{tabular}
\end{center}

Per \(n=45\) nodi il passo 2 produce quindi istogrammi \(15\times 15\), ossia
\(225\) coppie \((x,y,\hat{H})\) per grafo. Poiche' tutti i grafi hanno la stessa
dimensione, i pesi \(w_t\) di~\eqref{eq:sigl-step3} sono uniformi e la loss pesata
degenera in un MSE ordinario.

\subsection{Formato di uscita}

Il graphon stimato viene campionato dalla INR su una griglia regolare di celle
centrate,
\begin{equation}
    \widehat{W}_{ij} = f_{\theta^\star}\!\left(\frac{i+\tfrac12}{R}, \frac{j+\tfrac12}{R}\right),
    \qquad i,j = 0,\ldots,R-1,
\end{equation}
quindi sanificato secondo~\eqref{eq:sanitize} e salvato come immagine PNG in scala
di grigio a 8 bit, con nome \texttt{sample\_\{$\ell$\}.png}. Un file
\texttt{bootstrap\_metadata.json} registra, per ciascun graphon, il cluster di
origine, l'identificativo di replica e gli indici dei grafi effettivamente usati.

\begin{remark}[Quantizzazione implicita a 8 bit]
\label{rem:png}
Il passaggio attraverso il formato PNG introduce una quantizzazione non
documentata nella formulazione teorica: i valori vengono arrotondati a
\(\mathrm{round}(255\,\widehat{W}_{ij})/255\). Poiche' lo stadio 2 applica una
seconda quantizzazione a \(K \approx 10\) livelli, la risoluzione a 8 bit non e'
il fattore limitante; resta tuttavia una scelta implementativa da segnalare,
perche' rende il formato di scambio tra stadio 1 e stadio 2 lossy e non
esattamente invertibile.
\end{remark}

\section{Stadio 2, variante A: OGAD--D3PM}
\label{sec:d3pm-impl}

\subsection{Rappresentazione del dato}

Il graphon e' trattato come un'\emph{immagine categoriale} a canale singolo. Dato
\(\widehat{W}\in[0,1]^{R\times R}\), la quantizzazione e'
\begin{equation}
    \label{eq:d3pm-quant}
    b(w) = \mathrm{clip}\Bigl(\bigl\lfloor w\,(K-1) \bigr\rceil,\; 0,\; K-1\Bigr)
    \;\in\; \{0,\ldots,K-1\},
\end{equation}
dove \(\lfloor\cdot\rceil\) denota l'arrotondamento al pari intero piu' vicino. La
dequantizzazione e' esattamente l'inversa:
\begin{equation}
    \label{eq:d3pm-dequant}
    b^{-1}(k) = \frac{k}{K-1}.
\end{equation}
Si osservi che questa non e' una discretizzazione a bin di uguale ampiezza, ma un
arrotondamento al piu' vicino fra \(K\) \emph{livelli equispaziati} su \([0,1]\).
I livelli sono \(\{0, \tfrac{1}{K-1}, \ldots, 1\}\) e i confini impliciti sono i
punti medi \(\tfrac{2k+1}{2(K-1)}\); i due bin estremi hanno quindi ampiezza
dimezzata rispetto a quelli interni.

Il modello opera sull'intera matrice \(R\times R = 128\times128\): non vi e'
estrazione del triangolo superiore ne' mascheramento della diagonale. Le
\(16\,384\) celle sono trattate come variabili categoriali condizionalmente
indipendenti dato \(x_t\), e la simmetria e' un vincolo \emph{soft}, ereditato dai
dati e reimposto a posteriori tramite~\eqref{eq:sanitize}. Questo comporta uno
spreco di capacita' pari a circa un fattore 2, ed e' una delle direzioni di
miglioramento piu' immediate della pipeline.

\subsection{Processo diffusivo}

\begin{itemize}
    \item \textbf{Matrice di transizione}: uniforme, secondo~\eqref{eq:q-uniform},
    con \(K\) classi. La prior a \(t=T\) e' dunque uniforme su \(\{0,\ldots,K-1\}\),
    coerentemente con la distribuzione stazionaria del processo.
    \item \textbf{Schedule}: coseno secondo~\eqref{eq:cosine} con esponente
    \(p=1\), troncato superiormente:
    \[
        \beta_t = \min\!\left(1 - \frac{\bar\alpha(t/T)}{\bar\alpha((t-1)/T)},\; 0.999\right).
    \]
    \item \textbf{Passi}: \(T = 1000\), fissato nel codice e non esposto come
    iperparametro di configurazione.
    \item \textbf{Cumulate}: le \(\overline{Q}_t\) sono precalcolate esplicitamente
    come prodotti di matrici \(K\times K\) e memorizzate come buffer; il costo di
    memoria e' \(O(TK^2)\), trascurabile per \(K \leq 14\).
    \item \textbf{Campionamento}: sia in avanti che all'indietro il campionamento
    categoriale e' realizzato con il \emph{Gumbel-max trick},
    \(x = \argmax_i (\log \pi_i + g_i)\) con \(g_i \sim \text{Gumbel}(0,1)\).
    All'ultimo passo (\(t=1\)) il rumore di Gumbel viene azzerato, rendendo
    l'ultima transizione deterministica.
\end{itemize}

Durante l'addestramento i timestep sono estratti uniformemente in
\(\{1,\ldots,T-1\}\), quindi \(t=T\) non viene mai visitato.

\subsection{Rete di denoising}

Il denoiser e' un \emph{Diffusion Transformer} in stile DiT
\cite{peebles2023dit} con blocchi di tipo LLaMA \cite{touvron2023llama},
denominato \texttt{DiT\_Llama} nel repository \cite{d3pm_pytorch}:

\begin{center}
\begin{tabular}{ll}
    \toprule
    Componente & Configurazione \\
    \midrule
    Stem convoluzionale & 2 \(\times\) \texttt{Conv2d}(\(k=5\)) + SiLU + GroupNorm, \(1 \to d/2\) canali \\
    Patchify & patch \(2\times2\) \(\Rightarrow\) \(64^2 = 4096\) token \\
    Larghezza \(d\) & 512 \\
    Blocchi transformer & 5, con 16 teste \\
    Attenzione posizionale & RoPE \cite{su2024roformer} \\
    Feed-forward & SwiGLU \\
    Condizionamento su \(t\) & adaLN-Zero da embedding sinusoidale (256 frequenze) + MLP \\
    Testa finale & \texttt{Linear} inizializzata a zero, \(d \to 4\cdot 2K\) \\
    \bottomrule
\end{tabular}
\end{center}

L'input intero viene riscalato in \([-1,1]\) tramite
\(\tilde{x} = 2x/(K-1) - 1\), e l'uscita e' sommata a una versione \emph{gated}
della codifica one-hot dell'input:
\begin{equation}
    \label{eq:dit-skip}
    \mathrm{logits}(x_t)
    =
    \mathrm{nn}_\theta(x_t)
    + \mathrm{onehot}(x_t)\odot\bigl|1 + \gamma_\theta(x_t)\bigr|,
\end{equation}
cioe' la rete apprende un residuo attorno al denoiser identita'. Le shape sono:
input \((B,1,128,128)\) interi, output \((B,1,128,128,K)\) logit.

\subsection{Addestramento}

La loss e' quella ibrida~\eqref{eq:d3pm-hybrid}, con il termine variazionale
calcolato come media (non somma) delle KL su tutte le \(16\,384\) posizioni. La
griglia di iperparametri esplorata e':

\begin{center}
\begin{tabular}{ll}
    \toprule
    Iperparametro & Valori \\
    \midrule
    \texttt{lr} & \(\{2,3,4,5\}\times10^{-5}\) \\
    \texttt{num\_classes} \(K\) & \(\{10, 12, 14\}\) \\
    \texttt{hybrid\_loss\_coeff} \(\lambda\) & \(\{0.006, 0.01, 0.015\}\) \\
    \texttt{batch\_size} & 32 \\
    \texttt{dim} \(d\) & 512 \\
    \bottomrule
\end{tabular}
\end{center}

per un totale di \(4\times3\times3 = 36\) candidati per ogni configurazione SIGL.
L'ottimizzatore e' AdamW con clipping della norma del gradiente a \(5.0\) e mixed
precision; non e' usata alcuna media esponenziale dei pesi. Il protocollo di
early stopping monitora la loss di validazione (frazione \(0.2\), split
deterministico con seed \(0\)) con \texttt{patience} \(=25\),
\texttt{min\_delta} \(=5\times10^{-5}\) e al piu' \(300\) epoche; i pesi
dell'epoca migliore vengono ripristinati.

\begin{remark}
La loss di validazione e' essa stessa stocastica, perche' a ogni valutazione
vengono riestratti sia \(t\) sia il rumore di Gumbel. Il segnale di early stopping
e' quindi rumoroso, il che suggerisce di mediare su piu' passate o di fissare i
timestep di validazione come miglioria futura.
\end{remark}

\subsection{Generazione}

Il campionamento parte da rumore uniforme
\(x_T \sim \Uniform\{0,\ldots,K-1\}^{16\times1\times128\times128}\) e itera la
transizione inversa~\eqref{eq:x0-param} per \(t = T-1, \ldots, 1\), cioe' \(999\)
passi. Il tensore finale viene dequantizzato secondo~\eqref{eq:d3pm-dequant} e
sanificato secondo~\eqref{eq:sanitize}, producendo
\begin{center}
\texttt{data/predictions/graphons/tensors.pt} \quad di forma \([16,1,128,128]\).
\end{center}

\section{Stadio 2, variante B: OGAD--DiGress}
\label{sec:digress-impl}

\subsection{Riduzione del graphon a un grafo etichettato}

L'idea centrale di questa variante e' che un graphon quantizzato \emph{e'} un
grafo con archi categoriali. Formalmente si definisce l'immersione
\begin{equation}
    \label{eq:graphon-to-graph}
    \Phi_K : [0,1]^{R\times R} \longrightarrow \Graphs_{R,1,K},
    \qquad
    \Phi_K(\widehat{W}) = \bigl(X, E\bigr),
\end{equation}
dove:
\begin{itemize}
    \item il grafo sottostante e' il \emph{grafo completo} \(K_R\) su \(R=128\)
    nodi, privo di self-loop, con \(R(R-1) = 16\,256\) archi orientati;
    \item \(X = \one_{R\times 1}\), cioe' vi e' un \emph{unico} tipo di nodo
    (\(a=1\)) e i nodi non portano informazione;
    \item \(E_{ij} \in \{0,1\}^{K}\) e' la codifica one-hot della classe
    \begin{equation}
        \label{eq:digress-quant}
        \ell_{ij} = \bigl\lfloor \Pi(\widehat{W})_{ij}\,(K-1) \bigr\rceil
        \;\in\;\{0,\ldots,K-1\},
    \end{equation}
    con \(b = K\) tipi di arco;
    \item \(y\) e' vuoto (\(0\) feature globali).
\end{itemize}

La quantizzazione~\eqref{eq:digress-quant} e' formalmente identica
alla~\eqref{eq:d3pm-quant} di D3PM: entrambe le varianti condividono lo stesso
schema di discretizzazione, il che rende il confronto tra i due generatori
significativo. Con \(K=10\) i livelli rappresentabili sono
\(\{0, \tfrac19, \tfrac29, \ldots, 1\}\), ossia un passo di quantizzazione
\(\approx 0.111\).

\begin{remark}[Che cosa cambia rispetto a DiGress standard]
In DiGress applicato ai grafi molecolari, la classe di arco \(0\) rappresenta
\emph{l'assenza di arco} e le altre classi i tipi di legame; il grafo e' sparso.
Qui, al contrario, la classe \(0\) rappresenta la probabilita' \(W_{ij}=0\) ed e'
una classe come le altre: il grafo di supporto e' sempre completo. Di conseguenza
la distribuzione dei nodi e' degenere (\(n = 128\) con probabilita' 1) e il
conteggio delle classi d'arco viene calcolato includendo esplicitamente il bin
zero, invece di riservarlo ai non-archi come nell'implementazione originale.
\end{remark}

\subsection{Processo diffusivo}

Poiche' \(a=1\), la matrice \(Q_X^t\) e' l'identita' \(1\times1\) e il ramo dei
nodi e' inerte. Tutta la dinamica si svolge sul tensore degli archi, con
transizione \textbf{uniforme}:
\begin{equation}
    \label{eq:digress-QE}
    Q_E^t = (1-\beta_t)\,I_K + \frac{\beta_t}{K}\,\one\one^{\top},
    \qquad
    \overline{Q}_E^t = \bar\alpha_t\, I_K + (1-\bar\alpha_t)\,\frac{\one\one^{\top}}{K},
\end{equation}
con distribuzione limite uniforme sui \(K\) bin. Lo schedule e' coseno
secondo~\eqref{eq:cosine} con esponente \(p=2\) e \(T=1000\).

\begin{remark}[Uniforme anziche' marginale]
Questa e' una deviazione consapevole dal default di DiGress, che usa la
transizione marginale. Il codice della pipeline calcola comunque la distribuzione
marginale empirica dei bin (includendo il bin zero, come detto sopra), ma con la
configurazione \texttt{transition: uniform} tale quantita' non viene consumata.
La motivazione della scelta e' la comparabilita' con OGAD--D3PM, che usa
anch'esso una transizione uniforme; attivare la variante marginale richiede la
sola modifica del campo di configurazione ed e' la prima ablation naturale da
eseguire, dato che il Teorema 4.1 di DiGress ne prevede un vantaggio.
\end{remark}

Il rumore e' applicato al solo triangolo superiore di \(E\) e poi simmetrizzato,
quindi in questa variante --- a differenza di OGAD--D3PM --- la simmetria del
graphon e' garantita \emph{per costruzione} lungo tutta la traiettoria diffusiva.

\subsection{Rete di denoising e loss}

Viene usato il \texttt{GraphTransformer} di DiGress
\cite{vignac2023digress,dwivedi2021graphtransformer} senza alcuna modifica: 5
strati che aggiornano congiuntamente le feature di nodo \(X\), di arco \(E\) e
globali \(y\), con condizionamento reciproco tramite layer FiLM
\cite{perez2018film} e aggregazione PNA per le feature globali. Le dimensioni
effettive sono \(\mathrm{hidden\_mlp} = \{X{:}128, E{:}64, y{:}128\}\) e
\(\mathrm{hidden} = \{d_X{:}256, d_E{:}64, d_y{:}64,\ n_{\mathrm{head}}{:}8\}\).

Le \textbf{feature ausiliarie strutturali e spettrali sono disattivate}. La scelta
e' teoricamente motivata: il grafo di supporto e' una cricca su 128 nodi identica
per tutti i campioni, dunque conteggi di cicli, numero di componenti connesse e
spettro del laplaciano del supporto sono costanti e non veicolano alcuna
informazione. Le dimensioni di input della rete sono quindi \(X{:}1\), \(E{:}K\),
\(y{:}1\) (il singolo canale globale codifica il timestep normalizzato).

La loss e' la~\eqref{eq:digress-loss} con \(\lambda = 5\). Poiche' \(a=1\), la
softmax sui tipi di nodo e' degenere e \(\mathrm{CE}_X \equiv 0\); l'obiettivo
effettivamente ottimizzato e' dunque
\begin{equation}
    l(\hat{p}^G, G) = 5 \sum_{1\leq i,j\leq R} \mathrm{CE}\bigl(e_{ij}, \hat{p}^E_{ij}\bigr),
\end{equation}
ossia una pura cross-entropy sui \(K\) bin di probabilita'. Il criterio di
selezione del checkpoint e' invece la NLL variazionale di validazione
(\texttt{val/epoch\_NLL}), corrispondente al bound~\eqref{eq:d3pm-elbo}.

\subsection{Addestramento e generazione}

\begin{center}
\begin{tabular}{ll}
    \toprule
    Iperparametro & Valore \\
    \midrule
    \texttt{batch\_size} & 2 \\
    \texttt{accumulate\_grad\_batches} & 16 (batch efficace 32) \\
    \texttt{n\_epochs} & 200 \\
    \texttt{lr} & \(2\times10^{-4}\) \\
    ottimizzatore & AdamW, \texttt{amsgrad}, \(\text{wd} = 10^{-12}\) \\
    \texttt{n\_layers} & 5 \\
    \texttt{num\_bins} \(K\) & 10 \\
    \texttt{resolution} \(R\) & 128 \\
    campioni finali & 16 (100 nella configurazione estesa) \\
    EMA & disattivata \\
    \bottomrule
\end{tabular}
\end{center}

Lo split dei graphon e' \(80/10/10\) deterministico. Il campionamento parte da
rumore uniforme sui bin, gia' simmetrizzato con diagonale nulla, e collassa il
tensore finale con \(\argmax\) sui \(K\) canali. La post-elaborazione consiste in:

\begin{enumerate}
    \item \textbf{dequantizzazione} \(\ell \mapsto \ell/(K-1)\) --- si usa
    \emph{solo} il tensore degli archi \(E\), i tipi di nodo predetti sono
    scartati;
    \item \textbf{sanificazione} secondo~\eqref{eq:sanitize};
    \item \textbf{correzione di densita'} opzionale, attiva nella configurazione
    corrente: detta
    \begin{equation}
        \rho(W) = \frac{1}{R(R-1)}\sum_{i\neq j} W_{ij}
    \end{equation}
    la densita' fuori diagonale e \(\rho^\star\) la densita' media dei graphon di
    training prodotti da SIGL, ogni graphon generato viene riscalato come
    \begin{equation}
        \label{eq:density-match}
        W \;\longleftarrow\; \Pi\!\left(\frac{\rho^\star}{\rho(W)}\,W\right).
    \end{equation}
\end{enumerate}

\begin{remark}
La correzione~\eqref{eq:density-match} e' un intervento \emph{ad hoc} e va
riportato come tale: essa forza un momento del primo ordine del graphon generato a
coincidere con quello dei dati, ma non e' parte del modello probabilistico. Inoltre,
per via del clipping in \([0,1]\) contenuto in \(\Pi\), quando il fattore di scala
supera l'unita' la densita' ottenuta sottostima sistematicamente il bersaglio. La
pipeline registra le densita' prima e dopo la correzione in un file
\texttt{stats.json} per rendere l'effetto verificabile.
\end{remark}

L'uscita e' \texttt{data/predictions/graphons/digress/tensors.pt}, di forma
\([N,1,128,128]\), formato identico a quello prodotto dalla variante D3PM: i due
generatori sono quindi intercambiabili rispetto agli stadi 3 e 4.

\section{Confronto tra le due varianti}

\begin{center}
\begin{longtable}{p{0.28\textwidth}p{0.32\textwidth}p{0.32\textwidth}}
    \toprule
    & \textbf{OGAD--D3PM} & \textbf{OGAD--DiGress} \\
    \midrule
    \endhead
    Rappresentazione & immagine categoriale \(128\times128\) & grafo completo \(K_{128}\) con archi a \(K\) classi \\
    Quantizzazione & \(\lfloor w(K-1)\rceil\) & \(\lfloor w(K-1)\rceil\) (identica) \\
    \(K\) & \(\{10,12,14\}\) & \(10\) \\
    Entita' modellate & tutte le \(16\,384\) celle & triangolo superiore, \(8\,128\) archi \\
    Simmetria & soft, reimposta a posteriori & garantita per costruzione \\
    Transizione \(Q_t\) & uniforme & uniforme (marginale disponibile) \\
    Schedule & coseno, \(p=1\), \(T=1000\) & coseno, \(p=2\), \(T=1000\) \\
    Denoiser & DiT/LLaMA, 5 blocchi, \(d=512\), RoPE & GraphTransformer, 5 strati, \(d_X=256\), \(d_E=64\) \\
    Induttivo & localita' 2D via patch e convoluzioni & equivarianza a permutazione \\
    Loss & \(L_{\mathrm{vb}} + \lambda\,\mathrm{CE}\) & \(5\cdot\mathrm{CE}_E\) \\
    Feature ausiliarie & assenti & disattivate esplicitamente \\
    Selezione modello & griglia di 36 candidati + MMD & checkpoint a NLL minima \\
    Post-processing & sanificazione & sanificazione + density matching \\
    \bottomrule
\end{longtable}
\end{center}

La differenza concettuale piu' rilevante riguarda il \emph{bias induttivo}. Il
denoiser di D3PM tratta il graphon come un segnale bidimensionale e sfrutta la
localita' spaziale attraverso convoluzioni e patch: questo e' coerente con il
fatto che, dopo l'ordinamento SIGL, i graphon dei grafi Community presentano una
struttura a blocchi \emph{localmente regolare}. Il denoiser di DiGress e' invece
equivariante a permutazione, quindi indifferente all'ordinamento delle righe e
delle colonne; poiche' l'ordinamento e' esattamente l'informazione che SIGL ha
faticato a produrre, l'equivarianza rappresenta in questo contesto piu' un vincolo
che un vantaggio. Ci si attende quindi che la variante D3PM sia avvantaggiata sui
graphon ordinati, e questa e' un'ipotesi verificabile sperimentalmente.

Un secondo aspetto riguarda la scala dei dati. Con \(C=50\) e \(B=13\) SIGL
produce \(650\) graphon, di cui \(80\%\) usati per il training. La variante D3PM,
con batch \(32\), esegue circa \(16\) passi di ottimizzazione per epoca; la
variante DiGress, con batch \(2\) e accumulo \(16\), ne esegue circa \(16\)
altrettanti ma su un modello con costo per campione molto maggiore, dato che il
GraphTransformer deve gestire un tensore di archi con \(16\,256\) entita' per
grafo.

\section{Stadio 3: campionamento dei grafi}
\label{sec:sampling}

Sia
\(\mathcal{D}_{\mathrm{test}}=\{G^{\mathrm{test}}_1,\ldots,G^{\mathrm{test}}_n\}\)
il test set e \(m_i = |V(G^{\mathrm{test}}_i)|\) le dimensioni dei grafi
osservati. La distribuzione empirica delle dimensioni e'
\begin{equation}
    \hat{\mu}_M = \frac{1}{n}\sum_{i=1}^{n}\delta_{m_i}.
\end{equation}
Per generare il campione \(r\)-esimo, la pipeline seleziona ciclicamente un
graphon generato,
\begin{equation}
    \widetilde{W}^{(r)} = \widetilde{W}_{\,1 + ((r-1)\bmod S)},
\end{equation}
estrae \(M_r \sim \hat{\mu}_M\) con rigetto se \(M_r > R\), e campiona il grafo
secondo~\eqref{eq:graphon-sampling}--\eqref{eq:latent-index}:
\begin{equation}
    \label{eq:sampling-proc}
    \eta_1,\ldots,\eta_{M_r} \overset{\mathrm{i.i.d.}}{\sim} \Uniform([0,1]),
    \quad
    z_i = \lfloor R\eta_i\rfloor,
    \quad
    A^{(r)}_{ij} \sim \Bernoulli\Bigl(\Pi(\widetilde{W}^{(r)})_{z_i z_j}\Bigr)
    \;\; (i<j),
\end{equation}
simmetrizzando poi \(A^{(r)}\). Il dataset generato usato in valutazione e'
\(\widetilde{\mathcal{D}} = \{\widetilde{G}_1,\ldots,\widetilde{G}_B\}\) con
\(B = 100\) e seed fissato a \(0\).

Due dettagli meritano attenzione. Primo, la selezione \emph{ciclica} anziche'
casuale dei graphon garantisce un utilizzo bilanciato: con \(S=16\) graphon e
\(B=100\) grafi, i primi 4 graphon vengono usati 7 volte e i restanti 6.
Secondo, si usa un \emph{unico} vettore di variabili latenti \(\eta\) per entrambi
gli assi, il che e' la formulazione corretta di~\eqref{eq:graphon-sampling}:
campionare due vettori indipendenti produrrebbe una matrice non simmetrica e non
corrisponderebbe al modello \(W\)-random graph.

Poiche' nel dataset Community tutti i grafi di test hanno 45 nodi e \(R = 128\),
si ha \(M_r = 45\) deterministicamente e il rigetto non si attiva mai.

\section{Stadio 4: valutazione GGME}
\label{sec:ggme-impl}

\subsection{Descrittori}

Sono implementati tre descrittori, tutti calcolati in forma normalizzata:

\begin{center}
\begin{tabular}{lll}
    \toprule
    Descrittore & Binning & Normalizzazione \\
    \midrule
    \(\phi_{\mathrm{deg}}\): istogramma dei gradi & lunghezza variabile (\(\max\deg + 1\)) & somma unitaria (pmf) \\
    \(\phi_{\mathrm{clust}}\): coeff.\ di clustering & 200 bin su \([0,1]\) & integrale unitario \\
    \(\phi_{\mathrm{lap}}\): spettro laplaciano norm. & 200 bin su \([0,2]\) & integrale unitario \\
    \bottomrule
\end{tabular}
\end{center}

L'istogramma dei gradi ha lunghezza variabile e viene zero-padded alla lunghezza
massima prima del calcolo del kernel.

\begin{remark}[Scala delle metriche]
\label{rem:density}
I descrittori a binning fisso sono normalizzati come \emph{densita'} e non come
distribuzioni di probabilita': l'integrale, non la somma, e' unitario. Il fattore
di scala e' quindi \(200\) per il coefficiente di clustering (200 bin su un
intervallo di ampiezza 1) e \(100\) per lo spettro (200 bin su ampiezza 2), mentre
l'istogramma dei gradi e' una pmf ordinaria. Poiche' un kernel lineare e'
omogeneo di grado 2, i valori di MMD del clustering risultano scalati di un fattore
\(\approx 200^2 = 4\times10^4\) rispetto a quelli del grado. \textbf{Questo spiega
interamente} l'ordine di grandezza radicalmente diverso delle tre colonne nei
risultati (Capitolo~\ref{ch:risultati}) e implica che i valori non sono
confrontabili \emph{tra} descrittori, ma solo \emph{tra} modelli a parita' di
descrittore.
\end{remark}

\subsection{Protocollo di selezione}

Il protocollo GGME e' implementato come segue:

\begin{enumerate}
    \item \textbf{Perturbazione}: si applica la perturbazione \texttt{AddEdges},
    che aggiunge ogni coppia non adiacente in modo indipendente con probabilita'
    \(p\), a \(17\) livelli di intensita' \(p \in \{0.15, 0.20, \ldots, 0.95\}\)
    (seed \(42\)). Ogni livello produce una copia perturbata dei 100 grafi di test.
    \item \textbf{Griglia}: per ciascun descrittore si esplorano
    \(\{\text{gaussiano}, \text{lineare}\} \times \{\sigma = 0.1, 1, 10\}\).
    \item \textbf{Selezione}: si sceglie la combinazione che massimizza la
    correlazione di Pearson fra il livello di perturbazione \(p\) e la
    \(\widehat{\MMD}^2\) risultante.
    \item \textbf{Valutazione}: con la configurazione selezionata si calcola
    \(\widehat{\MMD}^2\) fra test set e ciascuno dei tre insiemi in confronto.
\end{enumerate}

Le configurazioni selezionate empiricamente sono riportate in
Tabella~\ref{tab:kernel-selection}.

\begin{table}[htbp]
\centering
\caption{Configurazioni di kernel selezionate dal protocollo di perturbazione.}
\label{tab:kernel-selection}
\begin{tabular}{lccr}
    \toprule
    Descrittore & Kernel & \(\sigma\) & Correlazione \\
    \midrule
    \(\phi_{\mathrm{deg}}\) & gaussiano & 10 & 0.8946 \\
    \(\phi_{\mathrm{clust}}\) & lineare & --- & 0.7638 \\
    \(\phi_{\mathrm{lap}}\) & lineare & --- & 0.7112 \\
    \bottomrule
\end{tabular}
\end{table}

\subsection{Selezione del candidato D3PM tramite MMD}
\label{sec:selection}

Nella variante D3PM, oltre alla selezione per loss di validazione, e' previsto un
secondo stadio di selezione basato sulle metriche. Per ciascuno dei 36 candidati
si generano 100 grafi dai rispettivi graphon e si calcolano le tre MMD rispetto
agli \(80\) grafi tenuti da parte nello split di Sezione~\ref{sec:sigl-impl}.
Il punteggio usato e' il rapporto medio rispetto a un bersaglio fisso:
\begin{equation}
    \label{eq:selection-score}
    \mathrm{score}(c)
    =
    \frac{1}{3}\sum_{d \in \{\mathrm{deg},\mathrm{clust},\mathrm{lap}\}}
    \frac{\widehat{\MMD}^2_d(c)}{\tau_d},
    \qquad
    \tau = (0.0125,\; 99.7936,\; 1.8530),
\end{equation}
dove i valori \(\tau_d\) sono le MMD ottenute da DiGress applicato direttamente ai
grafi. Il candidato con punteggio minimo viene promosso e i suoi graphon
sovrascrivono \texttt{tensors.pt}.

\begin{remark}
La normalizzazione per \(\tau_d\) e' necessaria proprio a causa
dell'Osservazione~\ref{rem:density}: senza di essa, la media aritmetica delle tre
MMD sarebbe interamente dominata dal termine di clustering. La scelta di usare
come normalizzatore i valori di DiGress rende il punteggio interpretabile
(\(\mathrm{score} < 1\) significa \enquote{mediamente meglio di DiGress}), ma
introduce anche una dipendenza da una esecuzione specifica del baseline, che
andrebbe idealmente sostituita da una normalizzazione intrinseca, ad esempio la
MMD fra training e test set raccomandata da GGME.
\end{remark}

\subsection{Termini di confronto}

Lo stadio GGME confronta tre insiemi di grafi con il test set:
\begin{itemize}
    \item \textbf{DiGress}: grafi generati da DiGress \emph{vanilla} applicato
    direttamente ai grafi Community. Questa variante usa 2 classi di arco
    (presenza/assenza), transizione marginale, 8 strati, feature ausiliarie
    strutturali e spettrali attive, e produce 160 campioni. Nel flusso
    orchestrato tramite script questo stadio viene saltato
    (\texttt{-\/-skip-digress}) e la colonna corrispondente risulta vuota;
    i valori riportati provengono da un'esecuzione dell'intero stack
    \texttt{docker compose}.
    \item \textbf{OGAD}: i 100 grafi campionati dai graphon generati secondo
    Sezione~\ref{sec:sampling}.
    \item \textbf{Paper}: i 400 grafi di riferimento distribuiti con GGME.
\end{itemize}

\begin{remark}
I tre insiemi non hanno la stessa cardinalita' (\(160\), \(100\), \(400\)). Lo
stimatore non distorto~\eqref{eq:mmd-unbiased} gestisce correttamente \(n\neq m\),
ma la varianza della stima differisce fra le tre colonne. La pipeline impone un
controllo di consistenza fra DiGress e OGAD, ma non fra questi e il riferimento
Paper: uniformare le cardinalita' e' una delle correzioni previste al protocollo.
\end{remark}

\section{Orchestrazione e riproducibilita'}

Ogni stadio e' incapsulato in un servizio Docker distinto, con immagini separate:
\texttt{ogad-prep-data} (SIGL), \texttt{ogad-train-model-b} (D3PM) oppure
\texttt{ogad-train-model-b-digress} (DiGress su graphon),
\texttt{ogad-train-model-a} (DiGress baseline) e \texttt{ogad-eval-metrics}
(GGME). Lo script \texttt{scripts/run\_ogad\_experiment.sh} esegue in sequenza:

\begin{algorithm}[htbp]
\caption{Esecuzione di un esperimento OGAD}
\begin{algorithmic}[1]
\State \(\texttt{RUN\_ID} \gets\) timestamp UTC; crea la directory di run
\State scrive il manifesto e copia la configurazione JSON in \texttt{runs/current/}
\State seleziona lo stack compose in base a \texttt{OGAD\_GRAPHON\_GENERATOR} \(\in\) \{\texttt{d3pm}, \texttt{digress}\}
\State \textbf{stadio 1} \Comment{servizio \texttt{prep\_data}}
\State \quad \texttt{python Main.py -{}-config \$CONFIG}
\State \textbf{stadio 2} \Comment{servizio \texttt{train\_model\_b}}
\If{generatore \(=\) \texttt{d3pm}}
    \State \quad \texttt{python ogad\_model\_finetune.py -{}-config \$CONFIG}
    \State \quad \texttt{python ogad\_model\_inference.py}
\Else
    \State \quad \texttt{python src/main.py +experiment=graphons \{override Hydra\}}
\EndIf
\State \textbf{stadi 3--4} \Comment{servizio \texttt{eval\_metrics}}
\State \quad \texttt{python main.py -{}-config \$CONFIG -{}-skip-digress}
\State scrive log e \texttt{metrics.jsonl} nella directory di run
\end{algorithmic}
\end{algorithm}

La configurazione sperimentale e' descritta da un unico file JSON i cui valori
sono liste; il blocco SIGL viene espanso in prodotto cartesiano generando run
indipendenti, mentre il blocco D3PM resta una griglia interna a ciascuna run. Un
driver parallelo distribuisce le run su piu' GPU con assegnazione round-robin.
La riproducibilita' e' garantita fissando i seed di \texttt{random}, \texttt{numpy}
e \texttt{torch} per ciascun candidato, e usando split deterministici basati su
seed espliciti in tutti i punti di partizionamento dei dati.

\chapter{Risultati preliminari}
\label{ch:risultati}

\section{Risultati riportati}

La Tabella~\ref{tab:results} riporta i valori di \(\widehat{\MMD}^2\) rispetto al
test set per la configurazione con \(C=200\) cluster, insieme a una seconda
esecuzione con \(C=50\) e bootstrap.

\begin{table}[htbp]
\centering
\caption{\(\widehat{\MMD}^2\) rispetto al test set. Valori piu' bassi indicano
maggiore somiglianza. Si veda l'Osservazione~\ref{rem:density} per la scala.}
\label{tab:results}
\begin{tabular}{lrrr}
    \toprule
    Descrittore & DiGress & OGAD & Paper \\
    \midrule
    \multicolumn{4}{l}{\emph{Configurazione \(C=200\), senza bootstrap}} \\
    \(\phi_{\mathrm{deg}}\) & 0.0125 & 0.6301 & 0.0034 \\
    \(\phi_{\mathrm{clust}}\) & 99.7936 & 2458.3862 & 16.0080 \\
    \(\phi_{\mathrm{lap}}\) & 1.8530 & 44.1202 & 0.5686 \\
    \midrule
    \multicolumn{4}{l}{\emph{Configurazione \(C=50\), con bootstrap}} \\
    \(\phi_{\mathrm{deg}}\) & --- & 0.1562 & 0.0034 \\
    \(\phi_{\mathrm{clust}}\) & --- & 215.876 & 16.008 \\
    \(\phi_{\mathrm{lap}}\) & --- & 2.8238 & 0.5686 \\
    \bottomrule
\end{tabular}
\end{table}

\section{Lettura dei risultati}

Il primo dato rilevante e' che l'introduzione del clustering con bootstrap
(\(C=50\), \(B\geq 12\)) migliora OGAD di un fattore compreso fra \(4\) e \(15\)
su tutti e tre i descrittori rispetto alla configurazione senza bootstrap. Questo
e' coerente con l'interpretazione della Sezione~\ref{sec:sigl-impl}: aumentare la
cardinalita' e la variabilita' del dataset di graphon migliora la stima di
\(p_\theta(W)\).

Il secondo dato e' che OGAD resta comunque distante dal riferimento Paper di
circa un ordine di grandezza su ciascun descrittore. Il divario e' particolarmente
marcato sul coefficiente di clustering, che misura la tendenza triadica. Questo e'
prevedibile dalla struttura del modello: il campionamento
\eqref{eq:sampling-proc} genera archi \emph{condizionatamente indipendenti} dati i
latenti, quindi il numero atteso di triangoli e' completamente determinato da
\(W\) e la varianza triadica dei grafi generati e' quella minima compatibile con il
modello \(W\)-random graph. Se il dataset reale contiene correlazioni di ordine
superiore non catturabili da un graphon, nessun miglioramento del generatore
diffusivo potra' colmarle: si tratterebbe di un limite del \emph{modello} e non
dell'\emph{algoritmo}.

Distinguere fra le due spiegazioni --- limite del graphon come modello, oppure
errore di stima e generazione --- e' l'esperimento diagnostico piu' importante da
eseguire. Il protocollo naturale e' calcolare la MMD fra il test set e i grafi
campionati direttamente dai graphon \emph{stimati} da SIGL, senza passare per il
generatore diffusivo. Tale valore costituisce un \emph{floor} raggiungibile:
se e' vicino a quello di OGAD, il collo di bottiglia e' lo stadio 1; se e'
molto piu' basso, lo e' lo stadio 2.

\section{Ipotesi sulle sorgenti di errore}

L'analisi implementativa del Capitolo~\ref{ch:pipeline} suggerisce un elenco
ordinato di cause candidate:

\begin{enumerate}
    \item \textbf{Risoluzione della quantizzazione.} Con \(K=10\) il passo di
    quantizzazione e' \(1/9 \approx 0.111\). Per un dataset la cui densita' media
    e' \(0.234\), questo significa che l'intero intervallo utile di probabilita' e'
    coperto da 2--3 livelli. La griglia esplorata (\(K \in \{10,12,14\}\)) e'
    probabilmente troppo grossolana: valori di \(K\) sensibilmente maggiori
    (32, 64) meritano di essere testati.
    \item \textbf{Prior uniforme.} Con transizione uniforme e \(K=10\), la
    distribuzione limite assegna probabilita' \(1/10\) a ciascun livello, ma la
    distribuzione empirica dei valori del graphon e' fortemente concentrata sui
    livelli bassi. Il modello deve quindi spendere gran parte della traiettoria
    inversa per correggere la prior. La transizione \emph{marginale} di DiGress e'
    progettata esattamente per questo caso ed e' gia' predisposta nel codice.
    \item \textbf{Numero di graphon indipendenti.} Anche con \(L=800\), i graphon
    provengono da \(C=50\) cluster: la diversita' effettiva e' limitata dal numero
    di cluster piu' che dal numero di repliche.
    \item \textbf{Doppia quantizzazione.} Il formato di scambio PNG a 8 bit
    (Osservazione~\ref{rem:png}) introduce una perdita di informazione evitabile,
    sostituibile con la serializzazione diretta dei tensori.
    \item \textbf{Modellazione ridondante in D3PM.} Modellare tutte le
    \(16\,384\) celle anziche' le \(8\,128\) del triangolo superiore raddoppia la
    dimensionalita' del problema senza aggiungere informazione.
\end{enumerate}

\chapter{Conclusioni provvisorie}

La pipeline OGAD propone una scomposizione modulare del problema di generazione
di grafi: prima si stimano graphon dai dati osservati, poi si apprende una
distribuzione sui graphon, infine si campionano grafi e li si valuta con metriche
strutturali. Questa impostazione rende esplicita la rappresentazione intermedia
appresa dal modello e permette di confrontare generatori discreti diversi nello
stesso protocollo sperimentale, come mostrato dalle due varianti D3PM e DiGress
che condividono lo stesso schema di quantizzazione, lo stesso stadio di
campionamento e lo stesso protocollo di valutazione.

I risultati preliminari indicano che la rappresentazione a graphon, nella forma
attualmente implementata, non e' ancora competitiva con la generazione diretta di
grafi. L'analisi dettagliata dell'implementazione suggerisce pero' che le cause
piu' probabili siano scelte di discretizzazione e di prior, entrambe correggibili
senza modifiche architetturali.

Il programma di lavoro che ne deriva e' il seguente:

\begin{itemize}
    \item \textbf{Diagnostica del collo di bottiglia}: misurare la MMD dei grafi
    campionati direttamente dai graphon stimati da SIGL, per separare l'errore di
    stima dall'errore di generazione;
    \item \textbf{Ablation sulla quantizzazione}: estendere la griglia di \(K\) a
    valori significativamente maggiori, verificando il trade-off fra risoluzione
    di probabilita' e difficolta' del problema categoriale;
    \item \textbf{Ablation sulla transizione}: confrontare transizione uniforme e
    marginale nella variante DiGress, dove la seconda e' gia' implementata;
    \item \textbf{Confronto equo fra le due varianti}: eseguire OGAD--D3PM e
    OGAD--DiGress sulla stessa configurazione SIGL, con lo stesso \(K\), lo stesso
    numero di graphon generati e lo stesso numero di grafi valutati;
    \item \textbf{Consolidamento del protocollo}: uniformare la cardinalita' dei
    tre insiemi confrontati, aggiungere la MMD training--test come scala di
    riferimento secondo la raccomandazione di GGME, ed estendere la valutazione
    alle perturbazioni diverse da \texttt{AddEdges};
    \item \textbf{Semplificazioni implementative}: eliminare il formato di scambio
    PNG e modellare il solo triangolo superiore nella variante D3PM.
\end{itemize}

\bibliographystyle{plain}
\bibliography{references}

\end{document}
